\chapter{Into the Browser}\label{ch:browser}

The final layer puts both networks on a static web page. There is no server: the game
rules, the networks and the rendering all run in the visitor's browser.

\section{Architecture}

\begin{center}
\begin{tikzpicture}[node distance=0.7cm and 1.1cm]
  \node[node soft, text width=2.4cm] (ckpt) {PyTorch\\checkpoint};
  \node[node box, right=of ckpt, text width=2.4cm] (onnx) {ONNX graph\\\footnotesize opset 17};
  \node[node fill, right=of onnx, text width=2.6cm] (ort) {ONNX Runtime Web\\\footnotesize WebAssembly};
  \node[node box, below=0.9cm of ort, text width=3.0cm] (rules) {\code{rules.js}\\\footnotesize rules, observations};
  \node[node box, left=of rules, text width=2.9cm] (loop) {\code{app.js}\\\footnotesize clock, input, panel};
  \node[node soft, left=of loop, text width=2.4cm] (canvas) {Canvas\\\footnotesize 60 frames/s};
  \draw[flow] (ckpt) -- node[lbl,above=1pt]{export} (onnx);
  \draw[flow] (onnx) -- (ort);
  \draw[flow] (rules) -- node[lbl,right]{observations} (ort);
  \draw[flow] (ort.south west) to[bend right=15] node[lbl,above left,pos=0.4]{moves} (loop.north east);
  \draw[flow] (loop) -- (rules);
  \draw[flow] (loop) -- (canvas);
\end{tikzpicture}
\end{center}

\begin{itemize}
\item \textbf{Export.} Each checkpoint is exported to ONNX (opset 17) with a dynamic batch
  dimension. ONNX Runtime then runs the export on random inputs, and the export
  script requires its logits and values to match PyTorch to within a relative $10^{-3}$.
\item \textbf{Inference.} ONNX Runtime Web 1.20.1 runs on WebAssembly with a single
  thread, which avoids the cross-origin isolation headers that multi-threading would
  require of the host. The duel network (8.2\,MB) loads with the page; the arena network
  (9.7\,MB) loads the first time Arena mode is opened.
\item \textbf{Rules.} \code{rules.js} is a single-game port of the engine, including the
  eight-plane and window observations, and is checked against the Python engine on
  6{,}400 transitions (Chapter~\ref{ch:verify}). What the page shows is the game the
  networks were trained on, not an approximation of it.
\end{itemize}

\section{Modes}

\begin{center}
\begin{tabular}{@{}llp{0.5\linewidth}@{}}
\toprule
\textbf{Mode} & \textbf{Board} & \textbf{Seats}\\
\midrule
AI vs AI & $11\times11$ & duel network in both seats, sampling from its policy\\
You vs AI & $11\times11$ & duel network (most likely move) against the player\\
AI vs bot & $11\times11$ & duel network against the greedy bot\\
Arena & $19\times19$ & 2--8 copies of the arena network, sampling; $f_{\min}$ equals the snake count\\
\bottomrule
\end{tabular}
\end{center}

\begin{warnbox}[Two deterministic copies mirror each other]
With both seats playing their most likely move, the first AI-vs-AI games ended in
identical head-on draws on turn four. The starts are symmetric, so two copies of one
deterministic policy make symmetric choices into the centre food. The all-AI modes
therefore \emph{sample} from the policy, as in training. The modes against a human or a
bot play the most likely move, as in evaluation.
\end{warnbox}

A side panel shows, live, the move probabilities of one network for its next turn and
its value estimate $V_\theta$, together with each snake's length and health.

\section{Continuous motion from a turn-based game}

The game advances in whole turns, but the page must look fluid, like a modern Snake
game. The fix is to separate logic from drawing: the rules move in discrete steps, and
the renderer interpolates between the last two states.

\begin{defbox}[Interpolated body]
Let $p_0,\dots,p_{L'-1}$ be a snake's cells before the latest turn and $c_0,\dots,c_{L-1}$
after it, so $c_i=p_{i-1}$ for $i\ge1$. At fraction $t\in[0,1]$ of the turn, the snake is
drawn as the polyline
\[
  \bigl(\,(1-t)\,p_0 + t\,c_0,\;\; c_1,\;\dots,\;c_{L-1},\;\; (1-t)\,p_{L'-1} + t\,c_{L-1}\,\bigr)
\]
through cell centres, stroked with round joins and caps.
\end{defbox}

At $t=1$ the head sits at $c_0$ and the tail at $c_{L-1}$. At the start of the next turn,
$t=0$, they sit at $p'_0 = c_0$ and $p'_{L-1}=c_{L-1}$: the drawing is continuous across
turn boundaries. When a snake has just eaten, the duplicated tail makes the last two
points coincide, so the tail correctly stays put.

\begin{warnbox}[Three bugs on the way to fluid]
\begin{enumerate}
\item \textbf{Ease-in, ease-out.} The first version eased each step with a smoothstep
  curve. Its velocity is zero at both ends of every turn, so the snake visibly stopped at
  every cell. Linear interpolation moves at constant speed.
\item \textbf{Stalls between turns.} Each turn waited for the network's next move before
  starting. The fix computes the AI's moves for the next turn while the current one
  animates, and drives the game from the animation loop on a fixed-cadence clock. A
  frame-by-frame measurement then showed the snake moving in every one of 90 consecutive
  frames.
\item \textbf{Diagonal tails.} The first polyline joined the moving tail straight to
  $c_{L-2}$, skipping $c_{L-1}$, the cell the tail is entering. At a corner that cut the
  corner diagonally. Including $c_{L-1}$, as in the definition above, keeps every segment
  on the grid.
\end{enumerate}
\end{warnbox}

Player input is buffered and read at the last moment before each turn. Arrow keys, WASD,
swipes, and clicks on the board, which steer toward the clicked cell, all feed the same
queue. A reversal onto the snake's own neck is ignored rather than allowed to kill it.
